% 数学ノート集・共通テンプレート v2（積分論2026年度第1部の版面を参考）
% A5判。章内で定義・定理・例・注意等の番号を共通化。
% コンパイル： platex -interaction=nonstopmode math_note_template_v2.tex (2回)
%              dvipdfmx math_note_template_v2.dvi
\documentclass[a5paper,10pt,onecolumn]{jsarticle}
\usepackage[dvipdfmx,top=22mm,bottom=20mm,inner=19mm,outer=18mm,headheight=13pt,headsep=6mm]{geometry}
\usepackage{amsmath,amssymb,amsthm,mathtools}
\usepackage{enumitem,fancyhdr,titlesec}
\usepackage[dvipdfmx,bookmarks=true,colorlinks=false]{hyperref}
\usepackage{pxjahyper}

% ----- 差し替える文書情報 -----
\newcommand{\NoteTitle}{Brown運動のマルチンゲール表現定理}
\newcommand{\NoteSeries}{数学ノート集}

% ----- 文字・組版 -----
% 参考PDFの質素な明朝＋Computer Modernの印象を引き継ぎながら，
% A5判で窮屈にならないよう本文の行間を少し確保する．
\linespread{1.13}
\setlength{\parindent}{1zw}
\setlength{\parskip}{0pt}
\setlength{\abovedisplayskip}{6pt plus 2pt minus 2pt}
\setlength{\belowdisplayskip}{6pt plus 2pt minus 2pt}
\setlength{\abovedisplayshortskip}{3pt plus 2pt minus 1pt}
\setlength{\belowdisplayshortskip}{4pt plus 2pt minus 1pt}
\setlist[enumerate]{topsep=3pt,itemsep=1pt,parsep=0pt,leftmargin=2.4zw}
\titleformat{\section}{\normalfont\normalsize\bfseries}{\thesection}{0.8zw}{}
\titlespacing*{\section}{0pt}{2.0ex plus .4ex}{.9ex}
\titleformat{\subsection}{\normalfont\normalsize\bfseries}{\thesubsection}{0.7zw}{}
\titlespacing*{\subsection}{0pt}{1.5ex plus .3ex}{.6ex}
\raggedbottom

% ----- 柱とページ番号：ページ上部、罫線なし -----
\fancypagestyle{mathnote}{%
  \fancyhf{}%
  \fancyhead[L]{\footnotesize \NoteSeries\ （\NoteTitle）}%
  \fancyhead[R]{\footnotesize \thepage}%
  \renewcommand{\headrulewidth}{0pt}%
  \renewcommand{\footrulewidth}{0pt}%
}
\pagestyle{mathnote}

% ----- 節内通し番号の定義・定理群 -----
\newtheoremstyle{mathnotetheorem}
 {4.5pt}{4.5pt}{\normalfont}{}%
 {\bfseries}{}{0.65zw}%
 {\thmname{#1}\thmnumber{ #2}\thmnote{（#3）}}
\theoremstyle{mathnotetheorem}
\newtheorem{theorem}{定理}[section]
\newtheorem{definition}[theorem]{定義}
\newtheorem{lemma}[theorem]{補題}
\newtheorem{proposition}[theorem]{命題}
\newtheorem{corollary}[theorem]{系}
\newtheorem{example}[theorem]{例}
\newtheorem{remark}[theorem]{注意}
\renewcommand{\proofname}{証明}
\renewcommand{\qedsymbol}{$\square$}
\makeatletter
\renewenvironment{proof}[1][\proofname]{%
  \par\pushQED{\qed}\normalfont\topsep4\p@\@plus2\p@\relax
  \trivlist\item[\hskip\labelsep\bfseries #1]\ignorespaces
}{\popQED\endtrivlist\@endpefalse}
\makeatother

% ----- 使うノートで適宜拡張する記号 -----
\newcommand{\R}{\mathbb{R}}
\newcommand{\norm}[1]{\lVert #1\rVert}
\newcommand{\abs}[1]{\lvert #1\rvert}

\begin{document}
\section*{読了後にできること}
本稿を読むための前提は，Brown運動のマルチンゲール表現，Itô公式．
本稿の具体的な到達課題は，具体的なマルチンゲールをItô積分で表現する．

\section{表現問題の設定}

$(B_t)_{0\le t\le T}$ を標準Brown運動とし，
$\mathcal F_t^B$ を
$\sigma(B_s:0\le s\le t)$ の通常の拡大とする．
このフィルトレーションは完備かつ右連続なものとする．
$L^2(\mathcal F_T^B)$ における内積は
$\langle X,Y\rangle=E[XY]$ とする．

\begin{definition}
$\mathcal H^2_T$ を
$\mathcal F_t^B$ に関する実数値予測可能過程 $H$ で
\[
 \norm H_{\mathcal H^2_T}^2
 =E\int_0^T H_s^2\,ds<\infty
\]
を満たすものの $dt\otimes P$ に関する同値類の空間とする．
It\^o積分
\[
 I(H)=\int_0^TH_s\,dB_s
\]
は $\mathcal H^2_T$ から $L^2(\mathcal F_T^B)$ への
線形等距離写像である．
\end{definition}

\begin{proposition}[It\^o積分の値域]\label{prop:range}
\[
 \mathcal R=\left\{\int_0^TH_s\,dB_s:
 H\in\mathcal H^2_T\right\}
\]
は $L^2(\mathcal F_T^B)$ の閉線形部分空間であり，
$\mathcal R\subset\{Z\in L^2(\mathcal F_T^B):EZ=0\}$ である．
\end{proposition}
\begin{proof}
線形性はIt\^o積分の線形性から従う．
$(I(H_n))$ が $L^2$ で収束するならCauchy列であり，
It\^o等長性により $(H_n)$ は
$\mathcal H^2_T$ のCauchy列である．
$\mathcal H^2_T$ は $L^2$ 空間なので完備であり，
ある $H$ に収束する．再びIt\^o等長性を用いれば
$I(H_n)\to I(H)$ となるので $\mathcal R$ は閉である．
単純過程のIt\^o積分の平均は零であり，
$L^2$ 極限で平均は連続に変化するため
任意の $H\in\mathcal H^2_T$ に対して $EI(H)=0$ となる．
\end{proof}

\section{Gaussian指数関数の稠密性}

\begin{definition}
決定論的な実数値階段関数 $h:[0,T]\to\mathbb R$ に対して
\[
 \mathcal E_t(h)=
 \exp\!\left(\int_0^t h(s)\,dB_s
   -\frac12\int_0^t h(s)^2\,ds\right)
\]
を\emph{確率指数関数}という．
\end{definition}

\begin{lemma}\label{lem:exponential}
決定論的な有界階段関数 $h$ に対して
\[
 \mathcal E_t(h)=1+\int_0^t h(s)\mathcal E_s(h)\,dB_s
\]
が成り立つ．特に
$\mathcal E_T(h)-1\in\mathcal R$ である．
\end{lemma}
\begin{proof}
$Y_t=\int_0^t h(s)\,dB_s$，
$v_t=\int_0^t h(s)^2\,ds$ と置く．
関数 $f(x,v)=\exp(x-v/2)$ に
多次元It\^oの公式を適用すると
\[
 d\mathcal E_t(h)=h(t)\mathcal E_t(h)\,dB_t
\]
となる．
実際，$dY_t=h(t)dB_t$，
$dv_t=h(t)^2dt$ なので，
二次変分の補正項 $h(t)^2dt/2$ が
$v_t$ からのドリフトを打ち消す．
またGaussian分布の指数モーメントにより
\[
 E[\mathcal E_t(h)^2]
 =\exp\!\left(\int_0^th(s)^2\,ds\right).
\]
従って
\[
 E\int_0^T h(t)^2\mathcal E_t(h)^2\,dt
 =\int_0^Th(t)^2
   e^{\int_0^t h(s)^2ds}\,dt<\infty.
\]
よって右辺の積分は二乗可積分なIt\^o積分であり，
$\mathcal E_T(h)-1\in\mathcal R$ を得る．
\end{proof}

\begin{lemma}[指数関数の稠密性]\label{lem:exp-dense}
次の線形部分空間は $L^2(\mathcal F_T^B)$ に稠密である：
\[
 \mathcal V=
 \operatorname{span}
 \left\{\mathcal E_T(h):
 h\text{ は決定論的な有界階段関数}\right\}.
\]
\end{lemma}
\begin{proof}
$\mathcal V^\perp=\{0\}$ を示す．
$Z\in L^2(\mathcal F_T^B)$ が
全ての $\mathcal E_T(h)$ に直交するとする．
有理数時刻（必要なら $T$ を加える）
$q_1,q_2,\dots$ を稠密に取り，
$\mathcal G_n=\sigma(B_{q_1},\dots,B_{q_n})$ と置く．
任意の $n$ と $u_1,\dots,u_n\in\mathbb R$ について
\[
 \sum_{j=1}^nu_jB_{q_j}
 =\int_0^T h_u(s)\,dB_s,\qquad
 h_u(s)=\sum_{j=1}^nu_j1_{(0,q_j]}(s)
\]
が成立する．従って直交性から
\[
 E\left[Z\exp\!\left(\sum_{j=1}^nu_jB_{q_j}\right)\right]=0
 \qquad(u\in\mathbb R^n)
 \tag{$*$}
\]
を得る．指数の決定論的な分散補正係数は正であるから，
それを取り除いても零の等式は保たれる．

$\mu_n$ を
$\mathbb R^n$ 上の有限符号付き測度
\[
 \mu_n(A)=E[Z1_{\{(B_{q_1},\dots,B_{q_n})\in A\}}]
\]
とする．$Z\in L^2$ とGaussianベクトルの指数モーメントを
Cauchy--Schwarzの不等式で評価すれば，
任意の $z\in\mathbb C^n$ に対して
$\int e^{z\cdot x}\,\mu_n(dx)$ は絶対収束し，
コンパクトな $z$ の集合上で一様に収束する冪級数展開を持つ．
従ってこの関数は各変数について整関数である．
式 $(*)$ から実数 $z=u\in\mathbb R^n$ で零なので，
一変数の正則関数の恒等定理を各変数に順次適用すると
全ての $z\in\mathbb C^n$ で零となる．
特に $z=i\xi$ とすると
$\mu_n$ のFourier変換が恒等的に零である．
有限符号付き測度のFourier変換の一意性より
$\mu_n=0$ である．
すなわち
$E[Z\mid\mathcal G_n]=0$ が成立する．

最後に
$\mathcal F_T^B$ は
$\sigma(\bigcup_n\mathcal G_n)$ の完備化と一致する．
これはBrown運動の標本路の連続性から，
任意の時刻 $t$ の $B_t$ が有理時刻の値の極限として
可測であることによる．
$\bigcup_nL^2(\mathcal G_n)$ は
$L^2(\mathcal F_T^B)$ に稠密である：
有限個の座標に関する有界単純関数の線形包の
$L^2$ 閉包は，
単調類定理と有界関数の近似によって
$L^2(\sigma(\bigcup_n\mathcal G_n))$ 全体となる．
従って $E[Z\mid\mathcal G_n]=0$ が全ての $n$ で
成立するなら，$Z$ は稠密部分空間に直交するので
$Z=0$ となる．
これで $\mathcal V^\perp=\{0\}$ を示した．
\end{proof}

\begin{remark}
証明で用いたGaussian指数関数の稠密性は，
Brown運動の自然フィルトレーションに固有の重要な性質である．
Fourier変換の一意性と複素解析の恒等定理を利用するのは，
有限次元Gaussian分布の指数関数系が完全であることを
示すためである．
\end{remark}

\section{マルチンゲール表現定理}

\begin{theorem}[Brown運動のマルチンゲール表現]\label{thm:mrt}
$Z\in L^2(\mathcal F_T^B)$ とする．
このときただ一つの $H\in\mathcal H^2_T$ が存在して
\[
 Z=E[Z]+\int_0^TH_s\,dB_s
 \qquad\text{ほとんど確実}
\]
が成り立つ．
従って
\[
 M_t=E[Z\mid\mathcal F_t^B]
\]
と置けば
\[
 M_t=E[Z]+\int_0^tH_s\,dB_s
 \qquad(0\le t\le T)
\]
が各固定時刻でほとんど確実に成り立ち，右辺は連続な修正を与える．
\end{theorem}
\begin{proof}
$\mathcal R$ を命題\ref{prop:range}の閉部分空間とし，
$\mathcal W=\mathbb R1+\mathcal R$ と置く．
$\mathcal R$ は平均零の確率変数から成るので
定数関数の空間と直交する．
従って $\mathcal W$ は閉線形部分空間である．
補題\ref{lem:exponential}によって
任意の決定論的有界階段関数 $h$ に対し
$\mathcal E_T(h)\in\mathcal W$ である．
補題\ref{lem:exp-dense}により
これらの指数関数の線形包は
$L^2(\mathcal F_T^B)$ に稠密であるから，
$\mathcal W=L^2(\mathcal F_T^B)$ となる．
従って任意の $Z$ は
$Z=c+\int_0^TH_s\,dB_s$ と表される．
期待値を取ると $c=EZ$ である．
二つの被積分過程 $H,G$ が同じ表現を与えるなら
It\^o等長性から
\[
 E\int_0^T|H_s-G_s|^2\,ds
 =E\left|\int_0^T(H_s-G_s)\,dB_s\right|^2=0
\]
なので $H=G$ が $dt\otimes P$ について成立する．
従って表現は一意である．

表現式の両辺に $\mathcal F_t^B$ に関する
条件付き期待値を取れば，
右辺のIt\^o積分過程がマルチンゲールであることから
\[
 E[Z\mid\mathcal F_t^B]
 =E[Z]+\int_0^tH_s\,dB_s
\]
を得る．右辺は連続な修正を持つ．なお，左辺の各時刻の版を任意に選ぶと，
区別不能性の意味で一致するとは限らない．
\end{proof}

\begin{corollary}[二乗可積分マルチンゲールの連続性]\label{cor:continuous}
$(M_t)_{0\le t\le T}$ を
Brown運動の自然フィルトレーションに関する
二乗可積分なマルチンゲールとする．
このとき $M$ は連続な修正を持つ．
\end{corollary}
\begin{proof}
マルチンゲール性により
$M_t=E[M_T\mid\mathcal F_t^B]$ である．
定理\ref{thm:mrt}を $Z=M_T$ に適用すると
$M_t=E[M_T]+\int_0^tH_s\,dB_s$ と表される．
右辺は連続な修正を持つから結論が従う．
\end{proof}

\section{具体的な表現}

\begin{example}[Brown運動の二次多項式]
$M_t=B_t^2-t$ と置くと
\[
 M_t=2\int_0^tB_s\,dB_s
\]
となる．
\end{example}
\begin{proof}
$f(x)=x^2$ にIt\^oの公式を適用すると
$d(B_t^2)=2B_t\,dB_t+dt$ である．
$B_0=0$ なので積分表示を得る．
また
$E\int_0^TB_s^2\,ds=\int_0^Ts\,ds=T^2/2<\infty$ だから
被積分過程は $\mathcal H^2_T$ に属する．
\end{proof}

\begin{example}[指数マルチンゲール]
$\theta\in\mathbb R$ とし
$M_t=\exp(\theta B_t-\theta^2t/2)$ と置くと
\[
 M_t=1+\int_0^t\theta M_s\,dB_s.
\]
\end{example}
\begin{proof}
補題\ref{lem:exponential}を
定数関数 $h(s)=\theta$ に適用すればよい．
\end{proof}

\begin{proposition}[二つの独立Brown運動の下での違い]
$B,W$ を独立な標準Brown運動とし，
$\mathcal G_t=\sigma(B_s,W_s:0\le s\le t)$ の通常の拡大とする．
$M_t=W_t$ は $\mathcal G_t$ に関する連続
二乗可積分マルチンゲールだが，
$W_T=\int_0^TH_s\,dB_s$ となる
$\mathcal G_t$-予測可能な $H\in\mathcal H^2_T$
は存在しない．
\end{proposition}
\begin{proof}
独立増分性から $W$ は $\mathcal G_t$ に関する
マルチンゲールである．
一方，有界単純 $\mathcal G_t$-予測可能過程
$H_s=\sum_j\xi_j1_{(t_j,t_{j+1}]}(s)$ に対して
\[
 E\left[W_T\int_0^TH_s\,dB_s\right]
 =\sum_j E[W_T\xi_j(B_{t_{j+1}}-B_{t_j})].
\]
$W_T=W_{t_j}+(W_T-W_{t_j})$ と分解する．
$\xi_jW_{t_j}$ は $\mathcal G_{t_j}$ 可測であり，
$B$ の増分は $\mathcal G_{t_j}$ と独立で平均零だから，
最初の項の期待値は零である．
さらに $(W_T-W_{t_j}, B_{t_{j+1}}-B_{t_j})$ は，
$\mathcal G_{t_j}$ と独立であり互いにも独立な
中心正規増分なので，
二番目の項の期待値も零である．
従って全体は $0$ となる．
単純過程の $L^2$ 稠密性とIt\^o等長性によって，
任意の $H\in\mathcal H^2_T$ について
$E[W_T\int_0^TH_s\,dB_s]=0$ となる．
もし $W_T=\int_0^TH_s\,dB_s$ なら
$E[W_T^2]=0$ となるが，$E[W_T^2]=T>0$ なので矛盾する．
\end{proof}

\begin{remark}
表現定理の仮定は，
積分に使うBrown運動とフィルトレーションが
正確に対応していることを要求する．
フィルトレーションに別の独立な雑音を追加すると，
その雑音によるマルチンゲールを
元のBrown運動だけで表現することはできない．
\end{remark}


\section{到達課題の完全解答}
\noindent\textbf{二次マルチンゲールの積分表示．}
\begin{proposition}
$M_t=B_t^2-t$ は二乗可積分マルチンゲールで，
\[
M_t=2\int_0^tB_s\,dB_s
\]
と表せる．
\end{proposition}
\begin{proof}
Itôの公式を $f(x)=x^2$ に適用すると
$d(B_t^2)=2B_t\,dB_t+dt$．
従って $dM_t=2B_t\,dB_t$ であり，
$M_0=0$ だから表示式が従う．
任意の有限時間 $T$ について
$E\int_0^T(2B_s)^2ds=4\int_0^Ts\,ds<\infty$
なので右辺は二乗可積分マルチンゲールである．
\end{proof}
\noindent\textbf{到達確認．}
$N_t=B_t$ をItô積分で表現せよ．
\par\noindent\textbf{解答．}
$B_t=\int_0^t1\,dB_s$ である．
\end{document}
